
\chapter{2017年上海卷（春）}

\section{填空题}

\begin{problem}
设集合 \(A=\{1,2,3\}\)，集合 \(B=\{3,4\}\)，则 \(A\cup B=\)\fillinblank{}.
\end{problem}

\begin{answer}
\(\{1,2,3,4\}\)
\end{answer}

\begin{problem}
不等式 \(\lvert x-1\rvert<3\) 的解集为\fillinblank{}.
\end{problem}

\begin{answer}
\((-2,4)\)
\end{answer}

\begin{problem}
若复数 \(z\) 满足 \(2\overline{z}-1=3+6\i\)（\(\i\) 是虚数单位），则 \(z=\)\fillinblank{}.
\end{problem}

\begin{answer}
\(2-3\i\)
\end{answer}

\begin{problem}
若 \(\cos\alpha=\dfrac13\)，则 \(\sin\left(\alpha-\dfrac{\pi}{2}\right)=\)\fillinblank{}.
\end{problem}

\begin{answer}
\(-\dfrac13\)
\end{answer}

\begin{problem}
若关于 \(x\)，\(y\) 的方程组 \(\begin{cases}
x+2y=4,\\
3x+ay=6
\end{cases}\)
无解，则实数 \(a=\)\fillinblank{}.
\end{problem}

\begin{answer}
\(6\)
\end{answer}

\begin{problem}
若等差数列 \(\{a_n\}\) 的前 \(5\) 项的和为 \(25\)，则 \(a_1+a_5=\)\fillinblank{}.
\end{problem}

\begin{answer}
\(10\)
\end{answer}

\begin{problem}
若 \(P\)，\(Q\) 是圆 \(x^2+y^2-2x+4y+4=0\) 上的动点，则 \(\lvert PQ\rvert\) 的最大值为\fillinblank{}.
\end{problem}

\begin{answer}
\(2\)
\end{answer}

\begin{problem}
已知数列 \(\{a_n\}\) 的通项公式为 \(a_n=3^n\)，则 \(\displaystyle\lim_{n\to\infty}\dfrac{a_1+a_2+\cdots+a_n}{a_n}=\)\fillinblank{}.
\end{problem}

\begin{answer}
\(\dfrac32\)
\end{answer}

\begin{problem}
若 \(\left(x+\dfrac{2}{x}\right)^n\) 的二项展开式的各项系数之和为 \(729\)，则该展开式中常数项的值为\fillinblank{}.
\end{problem}

\begin{answer}
\(160\)
\end{answer}

\begin{problem}
设椭圆 \(\dfrac{x^2}{2}+y^2=1\) 的左，右焦点分别为 \(F_1\)，\(F_2\)，点 \(P\) 在该椭圆上，则使得 \(\triangle F_1F_2P\) 是等腰三角形的点 \(P\) 的个数是\fillinblank{}.
\end{problem}

\begin{answer}
\(6\)
\end{answer}

\begin{problem}
设 \(a_1\)，\(a_2\)，\(\cdots\)，\(a_6\) 为 \(1\)，\(2\)，\(3\)，\(4\)，\(5\)，\(6\) 的一个排列，则满足 \(\lvert a_1-a_2\rvert+\lvert a_3-a_4\rvert+\lvert a_5-a_6\rvert=3\) 的不同排列的个数为\fillinblank{}.
\end{problem}

\begin{answer}
\(48\)
\end{answer}

\begin{problem}
设 \(a\)，\(b\in\R\)，若函数 \(f(x)=x+\dfrac{a}{x}+b\) 在区间 \((1,2)\) 上有两个不同的零点，则 \(f(1)\) 的取值范围为\fillinblank{}.
\end{problem}

\begin{answer}
\((0,3-2\sqrt2)\)
\end{answer}

\section{单选题}

\begin{problem}
函数 \(f(x)=(x-1)^2\) 的单调递增区间是（\quad）.

\choices
    {\([0,+\infty)\)}
    {\([1,+\infty)\)}
    {\((-\infty,0]\)}
    {\((-\infty,1]\)}
\end{problem}

\begin{answer}
D
\end{answer}

\begin{problem}
设 \(a\in\R\)，“\(a>0\)”是“\(\dfrac1a>0\)”的（\quad）条件.

\choices
    {充分非必要}
    {必要非充分}
    {充分}
    {既非充分也非必要}
\end{problem}

\begin{answer}
C
\end{answer}

\begin{problem}
过正方体中心的平面截正方体所得的截面中，不可能的图形是（\quad）.

\choices
    {三角形}
    {长方形}
    {对角线不相等的菱形}
    {六边形}
\end{problem}

\begin{answer}
A
\end{answer}

\begin{problem}
如图，正八边形 \(A_1A_2\cdots A_8\) 的边长为 \(2\)，若 \(P\) 为该正八边形边上的动点，则 \(\overrightarrow{A_1A_3}\cdot \overrightarrow{A_1P}\) 的取值范围为（\quad）.

\bitmapfigure[float=inline,max width=0.28\linewidth]{img/2017/shanghai_spring/q16_fig1.png}

\choices
    {\([0,8+6\sqrt2]\)}
    {\([-2\sqrt2,8+6\sqrt2]\)}
    {\([-8-6\sqrt2,2\sqrt2]\)}
    {\([-8-6\sqrt2,8+6\sqrt2]\)}
\end{problem}

\begin{answer}
B
\end{answer}

\section{解答题}

\begin{problem}
如图，长方体 \(ABCD-A_1B_1C_1D_1\) 中，\(AB=BC=2\)，\(AA_1=3\).

\begin{enumerate}
\item 求四棱锥 \(A_1-ABCD\) 的体积；
\item 求异面直线 \(A_1C\) 与 \(DD_1\) 所成角的大小.
\end{enumerate}

\bitmapfigure[max width=0.28\linewidth]{img/2017/shanghai_spring/q17_fig1.png}
\end{problem}

\begin{answer}
\begin{enumerate}
\item \(4\)
\item \(\arctan\dfrac{2\sqrt2}{3}\)
\end{enumerate}
\end{answer}

\begin{problem}
设 \(a\in\R\)，函数 \(f(x)=\dfrac{2^x+a}{2^x+1}\).

\begin{enumerate}
\item 求 \(a\) 的值，使得 \(f(x)\) 为奇函数；
\item 若 \(f(x)<\dfrac{a+2}{2}\) 对任意 \(x\in\R\) 成立，求 \(a\) 的取值范围.
\end{enumerate}
\end{problem}

\begin{answer}
\begin{enumerate}
\item \(a=-1\)
\item \([0,2]\)
\end{enumerate}
\end{answer}

\begin{problem}
某景区欲建造两条圆形观景步道 \(M_1\)，\(M_2\)（宽度忽略不计），如图，已知 \(AB\perp AC\)，\(AB=AC=AD=60\)（单位：米），要求圆 \(M_1\) 与 \(AB\)，\(AD\) 分别相切于点 \(B\)，\(D\)，圆 \(M_2\) 与 \(AC\)，\(AD\) 分别相切于点 \(C\)，\(D\).

\begin{enumerate}
\item 若 \(\angle BAD=60^\circ\)，求圆 \(M_1\)，\(M_2\) 的半径（结果精确到 \(0.1\) 米）；
\item 若步道 \(M_1\)，\(M_2\) 的每米造价分别为 \(0.8\text{ 千元}\) 和 \(0.9\text{ 千元}\)，如何设计圆 \(M_1\)，\(M_2\) 的大小，使总造价最低？最低总造价是多少？（结果精确到 \(0.1\text{ 千元}\)）.
\end{enumerate}

\bitmapfigure[max width=0.34\linewidth]{img/2017/shanghai_spring/q19_fig1.png}
\end{problem}

\begin{answer}
\begin{enumerate}
\item \(M_1\) 半径 \(34.6\text{ 米}\)，\(M_2\) 半径 \(16.1\text{ 米}\).
\item \(M_1\) 半径 \(30\text{ 米}\)，\(M_2\) 半径 \(20\text{ 米}\)，最低总造价为 \(42.0\text{ 千元}\).
\end{enumerate}
\end{answer}

\begin{problem}
已知双曲线 \(\Gamma:x^2-\dfrac{y^2}{b^2}=1\)（\(b>0\)），直线 \(l:y=kx+m\)（\(km\neq0\)），\(l\) 与 \(\Gamma\) 交于 \(P\)，\(Q\) 两点，\(P'\) 为 \(P\) 关于 \(y\) 轴的对称点，直线 \(P'Q\) 与 \(y\) 轴交于点 \(N(0,n)\).

\begin{enumerate}
\item 若点 \((2,0)\) 是 \(\Gamma\) 的一个焦点，求 \(\Gamma\) 的渐近线方程；
\item 若 \(b=1\)，点 \(P\) 的坐标为 \((-1,0)\)，且 \(\overrightarrow{NP'}=\dfrac32\,\overrightarrow{P'Q}\)，求 \(k\) 的值；
\item 若 \(m=2\)，求 \(n\) 关于 \(b\) 的表达式.
\end{enumerate}
\end{problem}

\begin{answer}
\begin{enumerate}
\item \(y=\pm\sqrt3x\)
\item \(k=\pm\dfrac12\)
\end{enumerate}
\end{answer}

\begin{problem}
已知函数 \(f(x)=\log_2\dfrac{1+x}{1-x}\).

\begin{enumerate}
\item 解方程 \(f(x)=1\)；
\item 设 \(x\in(-1,1)\)，\(a\in(1,+\infty)\)，证明：\(\dfrac{ax-1}{a-x}\in(-1,1)\)，且 \(f\!\left(\dfrac{ax-1}{a-x}\right)-f(x)=-f\!\left(\dfrac1a\right)\)；
\item 设数列 \(\{x_n\}\) 中，\(x_1\in(-1,1)\)，\(x_{n+1}=(-1)^{n+1}\dfrac{3x_n-1}{3-x_n}\)，\(n\in\N^*\)，求 \(x_1\) 的取值范围，使得 \(x_3\ge x_n\) 对任意 \(n\in\N^*\) 成立.
\end{enumerate}
\end{problem}

\begin{answer}
\begin{enumerate}
\item \(x=\dfrac13\)
\end{enumerate}
\end{answer}
